\documentclass[11pt,a4paper]{article}
\usepackage[T1]{fontenc}
\usepackage[utf8]{inputenc}
\usepackage[dvipsnames]{xcolor}
\usepackage{amsthm}
\usepackage{newtxtext,newtxmath}
\usepackage[margin=27mm,headheight=15pt]{geometry}
\usepackage{microtype,mathtools,booktabs,array,longtable}
\usepackage{enumitem,fancyhdr,titlesec,tcolorbox,listings}
\usepackage[colorlinks=true,linkcolor=MidnightBlue,citecolor=MidnightBlue,urlcolor=MidnightBlue]{hyperref}
\usepackage[nameinlink,noabbrev]{cleveref}
\definecolor{ink}{HTML}{172F48}
\definecolor{accent}{HTML}{176B78}
\definecolor{pale}{HTML}{F0F5F7}
\titleformat{\section}{\large\bfseries\color{ink}}{\thesection}{0.8em}{}
\titleformat{\subsection}{\normalsize\bfseries\color{ink}}{\thesubsection}{0.8em}{}
\pagestyle{fancy}
\fancyhf{}
\fancyhead[L]{\small\scshape A cyclic retract for D\.{z}yga's automata}
\fancyhead[R]{\small Research draft}
\fancyfoot[C]{\thepage}
\renewcommand{\headrulewidth}{0.3pt}
\setlength{\parskip}{4pt}
\setlength{\parindent}{14pt}
\setlength{\emergencystretch}{2em}
\setcounter{tocdepth}{2}
\numberwithin{equation}{section}
\newtheorem{theorem}{Theorem}[section]
\newtheorem{lemma}[theorem]{Lemma}
\newtheorem{proposition}[theorem]{Proposition}
\newtheorem{corollary}[theorem]{Corollary}
\newtheorem{conjecture}[theorem]{Conjecture}
\theoremstyle{definition}
\newtheorem{definition}[theorem]{Definition}
\newtheorem{example}[theorem]{Example}
\theoremstyle{remark}
\newtheorem{remark}[theorem]{Remark}
\newcommand{\Ak}{\mathcal A_k}
\newcommand{\Bk}{\mathcal B_k}
\newcommand{\Zr}{\mathbb Z/r\mathbb Z}
\newcommand{\rt}{\operatorname{rt}}
\newcommand{\rank}{\operatorname{rank}}
\newcommand{\dist}{\operatorname{dist}}
\newcommand{\ResetLang}{\operatorname{Reset}}
\newcommand{\id}{\operatorname{id}}
\newcommand{\eps}{\varepsilon}
\newcommand{\N}{\mathbb N}
\lstset{basicstyle=\ttfamily\small,breaklines=true,columns=fullflexible,
 frame=single,rulecolor=\color{black!15},backgroundcolor=\color{pale},
 keepspaces=true,showstringspaces=false}
\hypersetup{pdftitle={Synchronizing Dzyga's Automata by a Cyclic Retract},
 pdfauthor={Research draft prepared for Vladimir Reshetnikov},
 pdfsubject={An explicit universal reset construction and certified finite pair distances}}

\title{\vspace{-12mm}\color{ink}\textbf{Synchronizing D\.{z}yga's Automata\\by a Cyclic Retract}\\[5pt]
\large An explicit all-parameter reset construction,\\a nearby obstruction, and certified pair distances}
\author{Research draft prepared for Vladimir Reshetnikov}
\date{20 September 2026}

\begin{document}
\maketitle
\thispagestyle{empty}

\begin{abstract}
We study the binary automata $\Ak$ on $n=3k+6$ states that occur in
Conjecture~8 of Szyku\l a's 2026 survey of synchronizing automata. We give
a uniform proof that every $\Ak$, $k\geq1$, is synchronizing. The construction
uses the word $ab$: its unique recurrent component is a cycle of length
$r=k+4$, and $p=(ab)^r$ is an idempotent retraction onto this cycle.
Two projected letter actions become a reflection and a reflection with one
defect. Their composition produces an elementary merger, from which we
construct a reset word of length
\[
8k^2+54k+92=\frac89 n^2+\frac{22}{3}n+16.
\]
This resolves the synchronization clause for the figure-defined family,
not the conjectured exact compress-with-another threshold. We also prove
that shortening the tail by one state produces a nonsynchronizing family
with a dihedral permutation quotient. The known upper bound $4k+8$ for
compressing $q_0$ is reproved and attributed to prior work. Independent
integer certificates establish the exact threshold and all pair distances
for $1\leq k\leq200$. The universal matching lower bound remains unproved
here. Complete proofs, source-model conventions, executable checks, and
verification limitations are included.
\end{abstract}

\begin{tcolorbox}[colback=pale,colframe=accent,boxrule=0.6pt,arc=1mm,
 title={\textbf{Scope of the result}},fonttitle=\normalsize]
\textbf{Proved for every parameter:} a reset word for $\Ak$; a quadratic
reset-length upper bound; a permutation-quotient obstruction for the
one-state-shorter comparison family.\\[3pt]
\textbf{Not proved for every parameter:} that $4k+8$ is the \emph{minimum}
length merging $q_0$ with another state, or the proposed exact formula for
the pairwise merging diameter. The supplied finite certificates do not
replace these missing uniform lower-bound proofs.
\end{tcolorbox}

\vfill
\noindent\small\textbf{Research status.} This is a new research draft with an
explicit, independently checkable argument, not a peer-reviewed publication
or a proof-assistant formalization. Novelty is assessed against the sources
identified below, not asserted against all possible prior literature.
\normalsize
\newpage
\tableofcontents
\newpage

\section{The problem selected and the result obtained}
\label{sec:intro}

\subsection{A concrete open question in automata theory}

A deterministic finite automaton transforms its state each time it reads a
letter. A \emph{reset word} sends every possible initial state to the same
final state. The problem of finding such a word is therefore a question about
a language of words, not about choosing an initial or an accepting state.
Our automata are more precisely complete deterministic \emph{semiautomata};
the shorter conventional name will be used throughout.

Szyku\l a's survey \cite[Section~2.3, Figure~3 and Conjecture~8]{Szykula2026}
records a family $\Ak$ on $3k+6$ states, credited there to D\.{z}yga's 2018
master's thesis. Its conjecture combines universal synchronizability with an
exact shortest-word assertion for one distinguished state. We select the
first clause as the main research problem:
\begin{quote}
For the family $\Ak$ defined in \cref{sec:model}, does a reset word exist for
every integer $k\geq1$?
\end{quote}

The advantage of this target is that a positive answer can be made entirely
explicit. It is not necessary to characterize every shortest word, and it is
not necessary to infer a universal pattern from a finite search. One can
instead identify transformations induced by a few parameterized words and
then compose them to obtain a constant transformation.

An intervening online research manuscript \cite{MM2026} already proves the
all-parameter upper bound $4k+8$ for merging the distinguished state and
corrects the displayed witness word. It reports exact finite checks through
$k=60$, but explicitly leaves universal synchronizability open. Those upper
bounds and the correction are \emph{prior results}, not contributions claimed
by the present draft. We prove synchronization by a different structural
argument and extend the independently reproduced finite checks to $k=200$.

The original thesis has not been examined for this article. Thus historical
attribution to it is through the survey, and no assertion is made about the
precise wording or contents of that thesis. The available source record
supports treating the selected clause as open, but it cannot guarantee
absolute priority for the argument given here.

\subsection{Main theorem}

Words act on the right and are read from left to right: $q(uv)=(qu)v$.
A product of words below is a literal concatenation, even when we use the
same notation for the state transformations that the words induce.

\begin{theorem}[Explicit synchronization]\label{thm:main}
Let $k\geq1$ and let $\Ak$ be the complete binary automaton of
\cref{def:family}. Put
\begin{equation}\label{eq:mainwords}
 r=k+4,\qquad c=ab,\qquad p=c^r,\qquad
 e=a p b p c^{r-1}.
\end{equation}
Then the word
\begin{equation}\label{eq:resetword}
 R_k=p\bigl(e c^{r-1}\bigr)^{r-2}e
\end{equation}
sends every state to $q_2$. In particular, $\Ak$ is synchronizing and
\begin{equation}\label{eq:lengthmain}
 \rt(\Ak)\leq |R_k|=8r^2-10r+4
 =8k^2+54k+92
 =\frac89n^2+\frac{22}{3}n+16,
 \qquad n=3k+6.
\end{equation}
\end{theorem}

The notation $\rt(\Ak)$ denotes the length of a \emph{shortest} reset word.
The theorem supplies an upper bound, not a formula for that minimum. Its
purpose is an explicit synchronization proof. There is substantial room to
shorten the construction.

The proof has three exact steps. First, $p$ maps the whole state set onto a
cycle and fixes that cycle pointwise. Second, the words $ap$ and $bp$ act on
cycle indices as two almost matching reflections. Their discrepancy yields
an elementary merger $e$. Third, a rotation followed by this merger repeatedly
reduces the image until only one state remains. All three assertions are
proved below for an arbitrary parameter.

\section{Words, pair merging, and regular languages}
\label{sec:basics}

\subsection{Three different length parameters}

Let $\mathcal A=(Q,\Sigma,\delta)$ be finite and complete, with $|Q|\geq2$. For a nonempty
subset $S\subseteq Q$ and a word $w$, write
\[
 Sw=\{qw:q\in S\},\qquad \rank(w)=|Qw|.
\]
The empty word $\eps$ acts as the identity. For $x,y\in Q$ define
\begin{equation}\label{eq:distance}
 \dist(x,y)=\min\{\,|w|:xw=yw\,\},
\end{equation}
with value $\infty$ when the set is empty. The diagonal has distance zero.
For a distinguished state $q$, the compress-with-another threshold is
\begin{equation}\label{eq:tau}
 \tau(q)=\min_{x\in Q\setminus\{q\}}\dist(q,x).
\end{equation}
The reset threshold is instead
\[
 \rt(\mathcal A)=\min\{\,|w|:|Qw|=1\,\}.
\]
Finally, the pairwise merging diameter is
$D(\mathcal A)=\max_{x,y\in Q}\dist(x,y)$.

These parameters quantify different tasks. A word witnessing $\tau(q)$ may
merge just one pair and leave almost every other state distinct. A reset
word witnesses every pair's mergeability simultaneously. Whenever the
reset threshold is finite,
\begin{equation}\label{eq:thresholdorder}
 \tau(q)\leq D(\mathcal A)\leq\rt(\mathcal A).
\end{equation}
The first inequality follows from a minimum being at most a maximum over a
larger collection. The second follows because a reset word merges any chosen
pair. Neither inequality usually identifies the other quantities exactly.

\begin{lemma}[Pair criterion]\label[lemma]{lem:paircriterion}
A finite complete automaton is synchronizing if and only if every pair of
states can be merged by a word.
\end{lemma}
\begin{proof}
A reset word merges every pair. Conversely, begin with the image set $S=Q$.
If $S$ has more than one element, choose distinct $x,y\in S$ and a word $u$
with $xu=yu$. Applying $u$ to every member of $S$ cannot increase the image
size and identifies at least two members, so $|Su|<|S|$. Repeat on the new
image. Its positive integer size strictly decreases at each repetition,
so after at most $|Q|-1$ repetitions the concatenation of the chosen words
has singleton image.
\end{proof}

Our universal proof will not use finite pair enumeration. The lemma explains
why the independently supplied pair certificates also verify synchronization
on their finite parameter range.

\subsection{The reset language and a right-linear grammar}

The set
\[
 \ResetLang(\mathcal A)=\{w\in\Sigma^*:|Qw|=1\}
\]
is regular. Indeed, consider a deterministic automaton whose states are the
nonempty subsets of $Q$, whose initial state is $Q$, whose letter action is
$S\mapsto Sx$, and whose accepting states are the singleton subsets.
Induction on word length shows that its state after reading $w$ is exactly
$Qw$. This proves the claim without any minimization assumption.

It is also a two-sided ideal of the free monoid: if $w$ is a reset word and
$u,v$ are arbitrary words, then
\[
 Q(uwv)=((Qu)w)v\subseteq(Qw)v,
\]
and the right-hand side is a singleton. Completeness ensures the left-hand
side is nonempty, hence a singleton too.

For an explicitly grammatical description, introduce a nonterminal $X_S$
for each nonempty reachable subset, with start symbol $X_Q$. For each letter
$x$ include the production $X_S\to xX_{Sx}$, and include $X_S\to\eps$ exactly
when $S$ is a singleton. A derivation tracking its successive subscripts
shows that this right-linear grammar generates precisely
$\ResetLang(\mathcal A)$. Thus \cref{thm:main} gives a parameterized word in a
well-defined regular language. The grammar observation is elementary
background, not a new grammar-size theorem.

\section{The exact automaton model}
\label{sec:model}

\subsection{A total transition specification}

Fix $k\geq1$ and put
\begin{equation}\label{eq:parameters}
 n=3k+6,\quad M=2k+4,\quad N=3k+5,
 \qquad Q=\{q_0,q_1,\ldots,q_N\}.
\end{equation}
It is useful to call $q_0,\ldots,q_M$ the core, $q_M$ the junction, and
$q_{M+1},\ldots,q_N$ the tail. The tail has $k+1$ states.

\begin{definition}[The figure-complete family]\label[definition]{def:family}
The automaton $\Ak$ has alphabet $\{a,b\}$ and the following actions:
\begin{equation}\label{eq:a}
 q_i a=\begin{cases}
 q_0,&i=0,\\
 q_{i+1},&1\leq i<M,\ i\text{ odd},\\
 q_{i-1},&1\leq i<M,\ i\text{ even},\\
 q_{M+1},&i=M,\\
 q_3,&i=M+1,\\
 q_{\min(i+1,N)},&M+2\leq i\leq N,\ i\text{ even},\\
 q_{i-1},&M+2\leq i\leq N,\ i\text{ odd},
 \end{cases}
\end{equation}
\begin{equation}\label{eq:b}
 q_i b=\begin{cases}
 q_{i+1},&0\leq i<M,\ i\text{ even},\\
 q_{i-1},&0\leq i<M,\ i\text{ odd},\\
 q_2,&i=M,\\
 q_{\min(i+1,N)},&M+1\leq i\leq N,\ i\text{ odd},\\
 q_{i-1},&M+1\leq i\leq N,\ i\text{ even}.
 \end{cases}
\end{equation}
\end{definition}

Every pair $(q_i,x)$ is covered by exactly one clause of the appropriate
formula. The minimum in each outward tail rule is simply a compact way of
writing the endpoint self-loop. In particular,
\begin{equation}\label{eq:endpoint}
\begin{array}{c|cc}
 &q_Na&q_Nb\\\hline
 N\text{ even}&q_N&q_{N-1}\\
 N\text{ odd}&q_{N-1}&q_N
\end{array}
\qquad\text{and}\qquad q_0b=q_1.
\end{equation}

The survey's printed cases omit $q_0b$ and one endpoint transition; its
Figure~3 supplies them. Equations \eqref{eq:a}--\eqref{eq:endpoint} use that
figure-complete interpretation, also used in \cite{MM2026}. This convention
is essential: the theorem is not about an arbitrary completion of an
incomplete transition table. For example, changing $q_0b$ to $q_0$ would
make the single letter $b$ merge $q_0$ with $q_1$.

\subsection{Two immediate structural checks}

\begin{proposition}\label[proposition]{prop:basicmodel}
The automaton $\Ak$ is strongly connected. Both letter transformations have
rank $n-1$. The only distinct-state collisions of one letter are
\[
 \{q_4,q_{M+1}\}a=\{q_3\},\qquad
 \{q_3,q_M\}b=\{q_2\}.
\]
\end{proposition}
\begin{proof}
On $q_0,\ldots,q_{M-1}$, consecutive states are joined in both directions:
$b$ exchanges the pairs with even left endpoint, and $a$ exchanges the pairs
with odd left endpoint. Thus that segment is mutually reachable. From
$q_{M-1}$ an $a$ reaches $q_M$, and from $q_M$ a $b$ returns to $q_2$.
The transition $q_Ma=q_{M+1}$ reaches the tail; consecutive tail states
are again joined in both directions. Finally $q_{M+1}a=q_3$ returns to the
core. These routes prove strong connectivity.

The $b$-action consists of disjoint transpositions on the core segment,
disjoint tail transpositions and one possible endpoint fixed point, together
with $q_Mb=q_2$. The latter duplicates the image of $q_3$; every other image
has just one preimage, except the missing image $q_M$. Hence $b$ has rank
$n-1$ with the stated collision. For $a$, the exceptional chain is
\[
 q_{M-1}\longmapsto q_M\longmapsto q_{M+1}\longmapsto q_3.
\]
All remaining arrows are disjoint transpositions or fixed points. Its final
image $q_3$ is also the image of $q_4$, while $q_{M-1}$ is absent from the
image. This gives the other assertion.
\end{proof}

Strong connectivity does not itself prove synchronizability. Nor does the
fact that both letters have rank $n-1$. The comparison family in
\cref{sec:obstruction} will show how a nontrivial permutation quotient can
survive a very similar collection of arrows.

\section{The cyclic retract}
\label{sec:retract}

The main simplification appears after grouping letters into blocks $ab$.
Use the reparameterization
\begin{equation}\label{eq:rparameters}
 r=k+4\geq5,\qquad M=2r-4,\qquad N=3r-7,\qquad n=3r-6.
\end{equation}
Define $r$ distinct states
\begin{equation}\label{eq:cycle}
 s_0=q_0,\qquad s_j=q_{2j-1}\ (1\leq j\leq r-2),\qquad
 s_{r-1}=q_2,
 \qquad C=\{s_0,\ldots,s_{r-1}\}.
\end{equation}
Whenever a subscript on $s$ is interpreted modulo $r$, we say so or use an
expression already in $\Zr$.

\begin{lemma}[One recurrent cycle]\label[lemma]{lem:cycle}
For $c=ab$, the action on $C$ is
\begin{equation}\label{eq:rotation}
 s_jc=s_{j+1\bmod r}.
\end{equation}
Every state of $Q$ enters $C$ after at most $r-2$ applications of $c$.
The junction $q_M$ enters $C$ for the first time at $q_2$ after exactly
$r-2$ applications.
\end{lemma}
\begin{proof}
The beginning and end of the proposed cycle follow directly from the
transition rules:
\[
 q_0ab=q_1,\quad q_2ab=q_0,\quad q_{M-1}ab=q_2.
\]
For odd $i$ with $1\leq i\leq M-3$, the two letters send $q_i$ first to
$q_{i+1}$ and then to $q_{i+2}$. This proves the complete cycle
\[
 q_0\longmapsto q_1\longmapsto q_3\longmapsto\cdots
 \longmapsto q_{M-1}\longmapsto q_2\longmapsto q_0.
\]
The remaining even core states obey
\begin{equation}\label{eq:evencore}
 q_{2j}c=q_{2j-2}\qquad(2\leq j\leq r-3).
\end{equation}
Thus $q_{2j}$ reaches $q_2$ in $j-1\leq r-4$ blocks.

It remains to analyze the junction and tail. Let $E$ be the largest even
integer at most $N$, and $O$ the largest odd integer at most $N$. The
$c$-trajectory from the junction is
\begin{equation}\label{eq:tailwalk}
 \begin{split}
 q_M&\longmapsto q_{M+2}\longmapsto\cdots
       \longmapsto q_E\\
    &\longmapsto q_O\longmapsto\cdots
       \longmapsto q_{M+1}\longmapsto q_2.
 \end{split}
\end{equation}
The even portion lists the even tail states in increasing order, and the
odd portion lists the odd tail states in decreasing order; an endpoint is
listed only once when its portion has one state. Away from the endpoint,
the even tail indices increase by two under $c$
and the odd tail indices decrease by two. At the endpoint there are two
cases. If $N$ is even, $q_Na=q_N$ and $q_Nb=q_{N-1}$, so $q_Ec=q_O$.
If $N$ is odd, $E=N-1$ and $q_Ea=q_N$, $q_Nb=q_N$, again giving
$q_Ec=q_O$. At the exit, $q_{M+1}a=q_3$ and $q_3b=q_2$.
This proves every arrow in \eqref{eq:tailwalk}.

The displayed trajectory visits every one of the $N-M=r-3$ tail states
exactly once: first its even indices, then its odd indices. None belongs to
$C$. It therefore has $r-2$ edges from $q_M$ to its first point in $C$.
Every tail state is already on this trajectory and needs fewer edges.
Together with \eqref{eq:evencore}, this accounts for all of $Q$.
\end{proof}

The recurrent behavior has become completely regular despite the
parity-dependent endpoint. The parity changes only which of the last two
states is visited first in the folded tail trajectory.

\begin{lemma}[Idempotent retraction]\label[lemma]{lem:projection}
The word $p=c^r$ satisfies
\begin{equation}\label{eq:projection}
 Qp=C,\qquad s_jp=s_j\quad(0\leq j<r),\qquad p^2=p
 \quad\text{as transformations on }Q.
\end{equation}
Moreover,
\begin{align}
 q_{2j}p&=s_{r-j}&& (1\leq j\leq r-3),\label{eq:projecteven}\\
 q_Mp&=s_1,\label{eq:projectjunction}\\
 q_{M+1}p&=s_{r-2}.\label{eq:projectexit}
\end{align}
\end{lemma}
\begin{proof}
Every state enters $C$ by time $r-2$, hence certainly by time $r$.
On $C$, the word $c^r$ is one complete turn around a cycle of length $r$,
so it is the identity. These facts give both $Qp=C$ and $p^2=p$.

For \eqref{eq:projecteven}, $q_{2j}$ reaches $q_2=s_{r-1}$ after $j-1$
blocks, including the case $j=1$ in which it is there initially. The
remaining $r-j+1$ blocks put it at index
\[
 (r-1)+(r-j+1)\equiv r-j\pmod r.
\]
For the junction, \cref{lem:cycle} gives $q_Mc^{r-2}=q_2$. The two
remaining blocks send $q_2$ to $q_0$ and then $q_1=s_1$.
Finally $q_{M+1}c=q_2$; applying the remaining $r-1$ blocks gives index
$(r-1)+(r-1)\equiv r-2\pmod r$.
\end{proof}

It is important that $p$ fixes $C$ \emph{pointwise}; a word merely mapping
$Q$ into $C$ would not justify the same subsequent formulas. It is equally
important that $c^{r-1}$ is an inverse of $c$ only on $C$. The transformation
$c$ is not a permutation of the entire state set. No global inverse is used
anywhere in the construction.

\section{Two projected reflections and one elementary merger}
\label{sec:reflections}

A word $v$ mapping $C$ into itself induces a map on indices,
\[
 \overline v:\Zr\longrightarrow\Zr,
 \qquad s_jv=s_{\overline v(j)}.
\]
By \cref{lem:projection}, both $F=ap$ and $H=bp$ have this property, even
though $a$ and $b$ individually can leave $C$.

\begin{lemma}[Projected letters]\label[lemma]{lem:reflections}
The induced maps of $F=ap$ and $H=bp$ are
\begin{equation}\label{eq:FH}
 \overline F(j)=
 \begin{cases}
 1,&j=r-2,\\
 -j\pmod r,&j\neq r-2,
 \end{cases}
 \qquad
 \overline H(j)=1-j\pmod r.
\end{equation}
\end{lemma}
\begin{proof}
For $F$, $s_0a=q_0$ and $q_0p=s_0$, giving $\overline F(0)=0$.
For $1\leq j\leq r-3$, we have $s_j=q_{2j-1}$ and $s_ja=q_{2j}$;
\eqref{eq:projecteven} then gives $s_jap=s_{r-j}$.
At $j=r-2$, $s_j=q_{M-1}$, so $s_ja=q_M$ and
\eqref{eq:projectjunction} gives $s_jap=s_1$. This is the single exceptional
value. At $j=r-1$, $s_j=q_2$, $q_2a=q_1=s_1$, and $p$ fixes $s_1$.
Thus the formula for $F$ covers every index.

For $H$, the first two cases are $s_0b=q_1=s_1$ and $s_1b=q_0=s_0$.
For $2\leq j\leq r-2$,
\[
 s_jb=q_{2j-2}=q_{2(j-1)},\qquad
 q_{2(j-1)}p=s_{r-j+1},
\]
by \eqref{eq:projecteven}. Finally $s_{r-1}=q_2$ and $q_2b=q_3=s_2$.
These cases give $\overline H(j)=1-j$ modulo $r$ throughout.
\end{proof}

The map $\overline H$ is a genuine reflection of the cyclic indices.
The map $\overline F$ would be the reflection $j\mapsto-j$ but for one
value: index $r-2$ is sent to $1$ instead of $2$. This one-unit defect is
precisely what permits a reduction in rank.

\begin{lemma}[An elementary merger]\label[lemma]{lem:merger}
The word
\[
 e=F H c^{r-1}=a p b p c^{r-1}
\]
fixes every state of $C$ except $s_{r-2}$, and sends
$s_{r-2}$ to $s_{r-1}$. Equivalently,
\begin{equation}\label{eq:elementary}
 \overline e(j)=
 \begin{cases}
 r-1,&j=r-2,\\
 j,&j\neq r-2.
 \end{cases}
\end{equation}
\end{lemma}
\begin{proof}
For a nonexceptional index $j$, applying $\overline F$ and then
$\overline H$ gives $j\mapsto-j\mapsto1+j$. The final $c^{r-1}$ subtracts
one on $C$, so the result is $j$.
For $j=r-2$, the successive indices are
\[
 r-2\xmapsto{F}1\xmapsto{H}0\xmapsto{c^{r-1}}r-1.
\]
This is exactly \eqref{eq:elementary}.
\end{proof}

Although $F$ and $H$ are long words in the original alphabet, their
restrictions are very small algebraic objects. The use of a retraction
makes their compositions legitimate: after each projected action the
current state is back in $C$. There is no assumption that the same formulas
hold outside $C$.

\section{From a single merger to a reset word}
\label{sec:reset}

\subsection{A general cycle-sweeping lemma}

The final step is useful independently of the particular automaton.

\begin{lemma}[Cycle and adjacent merger]\label[lemma]{lem:sweep}
Let $C=\{s_0,\ldots,s_{r-1}\}$ with $r\geq3$. Suppose words $c,e$ map
$C$ into itself, $c$ advances indices by one modulo $r$, and $e$ fixes
all indices except an index $t$, which it sends to $t+1$ modulo $r$.
Put $h=e c^{r-1}$. Then
\begin{equation}\label{eq:sweep}
 \overline h(j)=
 \begin{cases}
 t,&j=t,\\
 j-1\pmod r,&j\neq t,
 \end{cases}
\end{equation}
and
\begin{equation}\label{eq:rankprofile}
 |Ch^m|=\max(1,r-m)\quad(m\geq0),\qquad
 Ch^{r-2}=\{s_t,s_{t+1}\},\qquad
 Ch^{r-2}e=\{s_{t+1}\}.
\end{equation}
If a word $p$ maps an entire state set $Q$ onto $C$, then
$p h^{r-2}e$ resets $Q$.
\end{lemma}
\begin{proof}
At index $t$, $e$ first advances to $t+1$, and $c^{r-1}$ returns to $t$.
Every other index is fixed by $e$ and then decremented by one. This proves
\eqref{eq:sweep}.

List the cyclic indices in the order
\[
 t-1,\ t-2,\ \ldots,\ t-(r-1)=t+1,\ t
 \qquad\text{in }\Zr.
\]
There are $r$ distinct entries. The transformation $h$ sends each entry to
the next and fixes the final entry. After $m$ iterations its image consists
of the final $r-m$ entries when $m<r$, and just the last entry when
$m\geq r-1$. This proves the rank formula and the description of
$Ch^{r-2}$. The map $e$ sends both of its remaining indices to $t+1$.
Prepending $p$ replaces the initial image $Q$ by $C$, proving the last
assertion.
\end{proof}

\subsection{Proof and length of the explicit construction}

\begin{proof}[Proof of \cref{thm:main}]
By \cref{lem:projection}, the first factor $p$ maps $Q$ onto $C$.
By \cref{lem:merger}, $e$ is the adjacent merger with exceptional index
$t=r-2$. Apply \cref{lem:sweep}: the suffix
$(e c^{r-1})^{r-2}e$ maps $C$ to $s_{r-1}=q_2$. Consequently
$QR_k=\{q_2\}$.

For the length calculation, all exponents in the word expression are
nonnegative integers and all products are literal concatenations. We have
\[
 |c|=2,\qquad |p|=2r,
\]
\[
 |e|=1+2r+1+2r+2(r-1)=6r,
 \qquad |e c^{r-1}|=8r-2.
\]
Hence
\begin{align*}
 |R_k|
 &=2r+(r-2)(8r-2)+6r\\
 &=8r^2-10r+4.
\end{align*}
Substituting $r=k+4$ gives $8k^2+54k+92$. Substituting
$r=n/3+2$ gives $8n^2/9+22n/3+16$. This proves every assertion of the
theorem.
\end{proof}

The state-set reduction is especially transparent:
\[
 Q\xmapsto{p}C
 \xmapsto{h^{r-2}}\{s_{r-2},s_{r-1}\}
 \xmapsto{e}\{s_{r-1}\}.
\]
The word has a short description even when its expanded length is large.
The implementation evaluates its transformations by repeated squaring,
using $O(n\log r)$ elementary state-array operations and $O(n)$ working
memory for the fixed number of arrays. Writing the expanded word instead
requires $\Theta(r^2)$ output symbols. These are different computational
tasks; the large-parameter checks use the compressed description.

\subsection{The smallest case worked out completely}

For $k=1$, the state set is $\{q_0,\ldots,q_8\}$ and $r=5$. The two
transition arrays, written as target indices, are
\begin{equation}\label{eq:k1arrays}
\begin{array}{c|rrrrrrrrr}
 i&0&1&2&3&4&5&6&7&8\\\hline
 a(i)&0&2&1&4&3&6&7&3&8\\
 b(i)&1&0&3&2&5&4&2&8&7
\end{array}
\end{equation}
The $c=ab$ cycle and its longest incoming trajectory are
\[
 0\to1\to3\to5\to2\to0,
 \qquad 6\to8\to7\to2.
\]
Thus the original state indices in the cycle order of $C$ are
$(0,1,3,5,2)$. On these five cycle positions,
not on the original state numbering, the relevant transformations are
\[
\begin{array}{c|rrrrr}
 j&0&1&2&3&4\\\hline
 \overline F(j)&0&4&3&1&1\\
 \overline H(j)&1&0&4&3&2\\
 \overline e(j)&0&1&2&4&4\\
 \overline h(j)&4&0&1&3&3
\end{array}
\]
The chain under $h$ is $2\to1\to0\to4\to3\to3$. Therefore
$Ch^3=\{s_4,s_3\}$ and $Ch^3e=\{s_4\}=\{q_2\}$. The word
$p h^3e$ has length $10+3\cdot38+30=154$. This example also illustrates why
mixing state labels with cycle indices would lead to incorrect formulas.

\section{Why one fewer tail state changes the answer}
\label{sec:obstruction}

A nearby nonsynchronizing example helps identify what the reset proof
actually exploits. It also gives a negative test for both the mathematics
and the implementation.

\begin{definition}[The shortened-tail family]
Keep $k\geq1$, $r=k+4$, and $M=2r-4$, but use the state set
\[
 Q'=\{q_0,\ldots,q_{N'}\},\qquad N'=3r-8=3k+4.
\]
Define $\Bk$ by the same formulas as \eqref{eq:a} and \eqref{eq:b}, with
$N'$ in place of $N$. In particular the new endpoint has the same
parity-based reflecting self-loop convention. The tail now has $k=r-4$
states, and $|Q'|=3k+5$.
\end{definition}

This operation is not simply restricting both old transformations to a
subset: the new endpoint's outward arrow is replaced by a self-loop.
For $k=1$ the tail consists of the single state $q_{M+1}$; its $a$-image
is still the exceptional exit $q_3$, and its $b$-image is itself.

\begin{theorem}[A dihedral permutation quotient]\label{thm:obstruction}
Every $\Bk$ is strongly connected but not synchronizing. More precisely,
there is a surjection $\pi:Q'\to\Zr$ such that
\begin{equation}\label{eq:quotient}
 \pi(q a)=-\pi(q),\qquad \pi(q b)=1-\pi(q)
 \qquad(q\in Q').
\end{equation}
\end{theorem}
\begin{proof}
Strong connectivity follows by the same core and tail routes as in
\cref{prop:basicmodel}; the one-state tail causes no exception to those
routes. Define the colors modulo $r$ by
\begin{equation}\label{eq:colours}
\begin{aligned}
 \pi(q_0)&=0,\\
 \pi(q_{2j-1})&=j &&(1\leq j\leq r-2),\\
 \pi(q_{2j})&=r-j &&(1\leq j\leq r-2),\\
 \pi(q_{M+2s+1})&=r-2-s &&(s\geq0,\ M+2s+1\leq N'),\\
 \pi(q_{M+2s})&=2+s &&(s\geq1,\ M+2s\leq N').
\end{aligned}
\end{equation}
These cases partition $Q'$. The states of $C$, which are all still present,
have colors $0,1,\ldots,r-1$, so $\pi$ is surjective.

On the ordinary core transpositions of $a$, the paired colors are $j$ and
$r-j$, which are negatives. The fixed state $q_0$ has color zero. On the
exceptional $a$-chain the colors are
\[
 r-2\longmapsto2\longmapsto r-2\longmapsto2,
\]
again changing sign at each step. A core $b$-pair with left endpoint
$q_{2j}$, $j\geq1$, has colors $r-j$ and $j+1$; these sum to $1$ modulo
$r$. The pair $q_0,q_1$ has colors $0,1$, and the junction transition
$q_Mb=q_2$ has colors $2,r-1$, again summing to $1$.

In the tail, an ordinary $a$-pair $q_{M+2s},q_{M+2s+1}$ has colors
$2+s,r-2-s$, which are negatives. An ordinary $b$-pair
$q_{M+2s+1},q_{M+2s+2}$ has colors $r-2-s,3+s$, which sum to $1$ modulo
$r$. The first tail state's exceptional $a$-exit has colors $r-2,2$,
as required. Thus only the endpoint self-loop remains to check.

If $r$ is even, $N'$ is even and $N'-M=r-4$ is even. The endpoint color
is $r/2$, fixed by $x\mapsto-x$, which is the letter $a$ that self-loops.
If $r$ is odd, $N'$ is odd and its color is $(r+1)/2$, fixed by
$x\mapsto1-x$, which is the self-looping letter $b$. The endpoint's other
letter is already covered by a tail pair, or by the exceptional exit when
$k=1$. This completes the verification of \eqref{eq:quotient}.

Both maps $x\mapsto-x$ and $x\mapsto1-x$ are permutations of $\Zr$.
By induction on word length, every word $w$ induces a permutation $g_w$
with $\pi(qw)=g_w(\pi(q))$. Consequently
\[
 \pi(Q'w)=g_w(\pi(Q'))=g_w(\Zr)=\Zr.
\]
In particular $|Q'w|\geq r>1$ for every word. No reset word exists.
\end{proof}

\begin{corollary}[A single endpoint defect]\label[corollary]{cor:defect}
Extend the same color formulas \eqref{eq:colours} to $\Ak$ by allowing its
one additional endpoint. The two relations \eqref{eq:quotient} hold at every
state-letter pair except the self-loop at $q_N$.
\end{corollary}
\begin{proof}
All non-self-loop verifications in the preceding proof remain valid.
If $r$ is even, then $N$ is odd, its color is $r/2$, and its $b$-loop would
need to send that color to $r/2+1$ instead of fixing it. If $r$ is odd,
then $N$ is even, its color is $(r+1)/2$, and its $a$-loop would need to send
that color to $(r-1)/2$. These colors differ modulo $r$ since $r\geq5$.
\end{proof}

This explains the usefulness of the projected reflections. In the shorter
family, a perfect reflection action survives as a quotient and forbids
synchronization. In the actual family, precisely one local arrow violates
that obstruction. The cyclic retraction turns the violation into an
explicit adjacent merger. A vague assertion that ``the letters almost act
as permutations'' would miss the decisive feature in either direction.

\section{The short merging word: a known result, reproved}
\label{sec:upper}

Let $\tau_k=\tau_{\Ak}(q_0)$. The word in the numerical part of the
original conjecture has one letter more than its stated length. The
corrected all-parameter upper bound is established in \cite{MM2026}; we
include a short proof using the trajectory already proved above.

\begin{proposition}[Known upper bound]\label[proposition]{prop:knownupper}
For every $k\geq1$, the word
\begin{equation}\label{eq:shortword}
 w_k=(ba)^{k+2}(ab)^{k+2}
\end{equation}
sends $q_0$ and $q_2$ to $q_2$. Thus
\begin{equation}\label{eq:tauupper}
 \tau_k\leq4k+8=\frac43n.
\end{equation}
\end{proposition}
\begin{proof}
Set $u=ba$, $c=ab$ and retain $r=k+4$. On the even core states,
\[
 q_{2j}u=q_{2j+2}\quad(0\leq j\leq r-3),\qquad q_Mu=q_1.
\]
Hence $q_0u^{r-2}=q_M$ and $q_2u^{r-2}=q_1$.
By \cref{lem:cycle}, $q_Mc^{r-2}=q_2$. The point $q_1=s_1$ is on the
$c$-cycle, so $q_1c^{r-2}=s_{r-1}=q_2$. Concatenating the blocks proves
both merging identities. Their combined length is $4(r-2)=4k+8$.
\end{proof}

The elementary free-monoid identity $(ab)^m a=a(ba)^m$ gives
\begin{equation}\label{eq:slip}
 (ab)^{k+2}a(ab)^{k+2}=a w_k,
 \qquad |a w_k|=4k+9.
\end{equation}
Since $q_0a=q_0$ and $q_1a=q_2$, this longer word merges $q_0$ with
$q_1$. The length discrepancy is not a counterexample to the numerical
bound $4k+8$; it is a distinction between a displayed witness, its partner,
and the minimum over all possible partners.

\begin{conjecture}[Remaining numerical clause]\label[conjecture]{conj:remaining}
For every $k\geq1$, $\tau_k=4k+8$.
\end{conjecture}

Only the lower inequality is missing after \cref{prop:knownupper}.
The universal reset word in \cref{thm:main} does not establish it: a word
that eventually merges \emph{all} states supplies no reason why a
\emph{particular pair} could not merge earlier than the proposed minimum.
The two questions must remain separate even though they originated in one
conjecture.

\section{Exact computations with independently checkable certificates}
\label{sec:certificates}

\subsection{A finite graph of unordered pairs}

For an $n$-state automaton, take all unordered pairs with repetition,
\[
 \mathcal P=\{\{q_i,q_j\}:0\leq i\leq j<n\}.
\]
There are $n(n+1)/2$ such vertices, including $n$ diagonal pairs. Each letter
induces one outgoing edge from each vertex. The merging distance of a pair
is exactly its directed distance to the diagonal set. This follows because
a path labeled by $w$ ends at $\{q_iw,q_jw\}$.

Reverse breadth-first search from all diagonal vertices computes these
distances in $O(n^2)$ time and space for a fixed two-letter alphabet: the
graph has $n(n+1)/2$ vertices and twice as many directed edges. The supplied
generator stores the reverse adjacency lists in integer arrays. It writes
one nonnegative integer for each pair. An unreachable pair would have value
$-1$ in the generator, but the verifier would reject such a certificate as
a certificate of complete mergeability.

\subsection{Why the certificate proves an exact distance}

\begin{lemma}[Bellman certificate]\label[lemma]{lem:certificate}
Let $\N$ include zero. Suppose a function $d:\mathcal P\to\N$ has value zero exactly on the
diagonal and satisfies, for every nondiagonal pair $P$,
\begin{equation}\label{eq:bellman}
 d(P)=1+\min\{d(Pa),d(Pb)\}.
\end{equation}
Then $d(P)$ equals the shortest merging length of $P$ for every $P$.
\end{lemma}
\begin{proof}
For either letter $x$, \eqref{eq:bellman} implies
$d(P)\leq1+d(Px)$. The same inequality holds on a diagonal because its
value and the values of both successors are zero. Along any word of length
$\ell$ that merges $P$, telescoping gives $d(P)\leq\ell+0$. Thus $d(P)$ is
a lower bound for every merging length.

If $P$ is nondiagonal, one of the two successor values is exactly $d(P)-1$.
Choose such a letter and repeat while the value is positive. A nonnegative
integer decreases by one at each step, so after exactly $d(P)$ steps it
reaches zero. Zero occurs only on the diagonal. The resulting word merges
$P$ and has length $d(P)$, proving the matching upper bound.
\end{proof}

This proof does not depend on how the values were discovered. In
particular the verifier need not rerun breadth-first search or trust the
generator's queue discipline. It checks only the finite table, the explicit
transition model, and the displayed local identity.

\subsection{Results actually verified}

For every integer $1\leq k\leq200$, the supplied certificates verify
\begin{equation}\label{eq:finitechecks}
 \tau_k=4k+8,\qquad \dist(q_0,q_1)=4k+9,
\end{equation}
and identify the complete set of optimal partners of $q_0$ as
$\{q_2,q_{2k+5}\}$. They also verify, on the same finite range,
\begin{equation}\label{eq:diameter}
 D(\Ak)=
 \begin{cases}
 (5k^2+32k+47)/4,&k\text{ odd},\\
 (5k^2+34k+44)/4,&k\text{ even}.
 \end{cases}
\end{equation}
The prior note \cite{MM2026} reports the threshold checks through $k=60$
and the diameter formula through $k=16$. Here all values are independently
recomputed and checked, not copied from that source.

\begin{table}[htbp]
\centering
\small
\begin{tabular}{rrrrrr}
\toprule
$k$&$n$&$\tau_k$&$\dist(q_0,q_1)$&$D(\Ak)$&$|R_k|$\\
\midrule
1&9&12&13&21&154\\
2&12&16&17&33&232\\
3&15&20&21&47&326\\
4&18&24&25&65&436\\
5&21&28&29&83&562\\
10&36&48&49&221&1432\\
20&66&88&89&681&4372\\
40&126&168&169&2351&15052\\
60&186&248&249&5021&32132\\
100&306&408&409&13361&85492\\
150&456&608&609&29411&188192\\
200&606&808&809&51711&330892\\
\bottomrule
\end{tabular}
\caption{Selected certified pair-distance results. The last column is the
length of the universal construction, not a computed shortest reset length.
The archive contains every intervening parameter, not only these rows.}
\label{tab:results}
\end{table}

The largest automaton in the exhaustive pair calculation has $606$ states
and $183{,}921$ pair vertices. Across the $200$ certificates there are
$12{,}486{,}300$ pair entries, each checked against either the diagonal
condition or the Bellman identity. All entries are finite and all the stated
identities passed.

Separately, every structural identity in the universal construction was
checked for $k=1,\ldots,300$ and at $k=500,1000,10000$. Those tests include
the retraction image, the two projected maps, the elementary merger, selected
ranks of powers of the sweeping map, and the final singleton image. They
also check the shortened-family quotient and the unique endpoint defect.
For $k=1,\ldots,10$, the reset words were expanded and evaluated both as
compositions of state arrays and by direct state-by-state simulation.

These computations support implementation fidelity and help catch indexing
errors. The proofs in \cref{sec:retract,sec:reflections,sec:reset,sec:obstruction}
are the justification for universal statements. Testing the word at
$k=10000$ is not being presented as a substitute for induction or a symbolic
argument.

\subsection{Implementation independence and negative controls}

The generator and verifier build the transition functions differently.
The generator uses the piecewise rules. The verifier initializes identity
maps, inserts disjoint transpositions, and then inserts the exceptional
junction arrows. It does not import the generator or its constructor.
The structural test program compares the two constructions directly.

The verifier checks every entry using ordinary Python integers and explicit
exceptions. Its safety conditions are not written as \texttt{assert}
statements that disappear under \texttt{python -O}. It also rejects three
deliberately corrupted versions of the smallest certificate: a nonzero
diagonal, an incremented nondiagonal distance, and a nondiagonal distance
replaced by zero. Random-word tests with a fixed seed check the
left-to-right composition convention against a second evaluator.

This is an independent executable checker in the algorithmic sense, not a
separately developed or proof-assistant-verified implementation. A shared
misreading of the source figure would still be a possible error unless the
source-model correspondence were checked; this is why the complete
transition specification and endpoint discussion are included in the
article itself.

\section{What the argument does not yet prove}
\label{sec:remaining}

\subsection{The universal lower bound}

To finish \cref{conj:remaining}, it would suffice to produce, for arbitrary
$k$, a nonnegative potential $L_k$ on unordered pairs such that it vanishes
on the diagonal and satisfies
\begin{equation}\label{eq:potential}
 L_k(P)\leq1+L_k(Pa),\qquad
 L_k(P)\leq1+L_k(Pb),
\end{equation}
together with
\begin{equation}\label{eq:potentialstart}
 L_k(\{q_0,q_j\})\geq4k+8\qquad(1\leq j\leq N).
\end{equation}
The telescoping half of \cref{lem:certificate} would then imply the missing
lower bound. Equality in a Bellman equation would not be necessary for
this purpose. A piecewise formula proving only the lower bound might be
simpler than a formula for all exact distances.

The finite distance tables provide examples of such potentials for the
checked parameters, but this article does not extract a universally valid
formula from them. Nor does it give a finite parameter-independent partition
of all pair configurations with verified affine transition inequalities.
Those are precise missing steps, not consequences of the reset construction.

\subsection{Why the synchronization proof cannot be reversed}

The retraction $p$ deliberately forgets transient position. Many states
outside $C$ acquire the same projected state, and that is useful for
constructing a reset word. For a shortest-word lower bound, however,
forgetting transient position can erase the very time cost one needs to
measure. Showing that a particular long word collapses a projected cycle
does not imply that every shorter word fails to merge the distinguished
state with some partner.

Likewise, the dihedral coloring in \cref{sec:obstruction} is an exact
obstruction only for the shortened family. In $\Ak$ its one defective
self-loop changes the color relation. A lower-bound argument based on that
coloring would have to track both when a trajectory reaches the defect and
what transient information remains after leaving it. Treating all letters
as reflections throughout a proposed short word would be false precisely
at that critical transition.

A further source of misleading lower bounds is symmetrizing the pair graph.
Allowing a transition to be traversed backward can only shorten distances,
so undirected distance is a legitimate relaxation but need not be sharp.
One must not replace a directed shortest-path problem by a symmetric
geometric picture and then silently retain the claimed exact length.

\subsection{Other conjectural patterns and distinct neighboring questions}

The formula \eqref{eq:diameter} is certified here only through $k=200$.
Its algebraic leading term would be $(5/36)n^2$ after substituting
$n=3k+6$, but finite agreement does not establish an asymptotic theorem for
$D(\Ak)$. The same limitation applies to the observed optimal-partner
classification outside the certified range.

Triple rendezvous time asks for a shortest word merging \emph{some three}
initial states. A recent preprint by Zhang \cite{Zhang2026} gives exact
$\lfloor4n/3\rfloor$ values for a different family. This is a neighboring
invariant, not the distinguished-state threshold of $\Ak$, and not the
universal reset problem solved by our construction. Coinciding numerical
expressions do not identify the automata or the quantified tasks.

Finally, the quadratic bound \eqref{eq:lengthmain} is not presented as
optimal. The construction spends many letters restoring states to the
cyclic retract after each projected action. Shorter projected words or
more economical sweeps could lower it. The exact reset threshold itself
has not been computed by the pair-distance procedure: it would require
an appropriate subset calculation or another proof.

\section{Conclusion}

The selected open synchronization question has an explicit answer for the
figure-complete family: every $\Ak$ resets under
\[
 \boxed{\;R_k=p(e c^{r-1})^{r-2}e,\quad
 c=ab,\quad p=c^r,\quad e=a p b p c^{r-1},\quad r=k+4.\;}
\]
The mechanism is a cyclic retraction, two projected reflections differing
at one index, and an elementary adjacent merger. It is uniform in the
parameter and includes both endpoint parities without an appeal to
experimental extrapolation.

The shortened-tail family supplies a matching structural warning: restoring
a perfect dihedral quotient prevents synchronization altogether. This
comparison identifies the endpoint defect rather than strong connectivity
or near-involutivity as the feature used by the proof.

The larger numerical conjecture remains a separate problem. Its upper
bound is prior work and is reproved here; its lower bound is independently
certified through $606$ states, but not proved for arbitrary $k$. The
archive is intended to make both the affirmative theorem and that boundary
of knowledge easy to inspect and reproduce.

\appendix
\section{Reproducing the artifacts}
\label{app:reproduce}

All Python programs use only the standard library and require Python~3.10
or newer. From the archive's top directory, run
\begin{lstlisting}
python code/test_structure.py
python code/verify_certificates.py --max-k 200
\end{lstlisting}
The second command checks the existing certificates without regenerating
them. To reproduce the data and CSV summary, run
\begin{lstlisting}
python code/generate_certificates.py --max-k 200
python code/verify_certificates.py --max-k 200
\end{lstlisting}
On a constrained execution service the generator may instead be split into
two runs:
\begin{lstlisting}
python code/generate_certificates.py --max-k 150
python code/generate_certificates.py --min-k 151 --max-k 200
\end{lstlisting}
The first run replaces the summary and the second appends to it. Do not
append the same range twice without regenerating the summary from $k=1$.

The certificate for parameter $k$ is
\texttt{certificates/k\_XXXX.json.gz}, where \texttt{XXXX} is $k$ padded to
four decimal digits. Its distance array uses index
\[
 I(i,j)=\frac{\max(i,j)(\max(i,j)+1)}2+\min(i,j).
\]
Thus pairs are listed by increasing larger endpoint and then by increasing
smaller endpoint. The verifier derives transitions independently from $k$;
transition arrays supplied by the generator are not trusted or needed.

The file \texttt{results/summary.csv} contains one row for each parameter.
The three text logs record data generation, certificate checking, and
structural testing. These logs report completed runs, not predicted results.
The directory \texttt{code/} includes a guarded expansion routine for small
reset words and a transformation-based evaluation routine for large ones.
The file \texttt{PROOF\_AUDIT.md} maps every main claim to its proof,
computational status, and relevant limitation.

To rebuild the document with a standard TeX Live installation:
\begin{lstlisting}
pdflatex -halt-on-error article.tex
pdflatex -halt-on-error article.tex
pdflatex -halt-on-error article.tex
\end{lstlisting}
The bibliography is embedded below, so no bibliography database or network
access is needed to compile. A separate \texttt{references.bib} is supplied
for reuse, not required by this build.

\clearpage
\section{A compact algebraic audit}

For checking the proof without following the discussion in its original
order, the exact dependencies are as follows. The transition formulas imply
the cycle and tail trajectory in \cref{lem:cycle}. Their lengths imply
$p=c^r$ is an idempotent retraction, together with the three projection
identities of \cref{lem:projection}. Those identities alone determine the
actions of $ap$ and $bp$ in \cref{lem:reflections}. Composing those two
maps and the inverse cycle rotation on $C$ yields the one-point merger of
\cref{lem:merger}. The path argument in \cref{lem:sweep} then proves the
reset identity. No computational assertion is an input to this chain.

There are three conventions on which the argument is especially sensitive.
First, the letter $a$ is applied \emph{before} $b$ in $c=ab$. Second,
$q_0b=q_1$ and the nonloop endpoint arrow returns to the previous state.
Third, indices modulo $r$ label the ordered states $s_j$ of the retract;
they are not the original labels on $q_i$. Each convention is exposed in
both the formulas and the executable examples.

\begin{thebibliography}{9}

\bibitem{Szykula2026}
Marek Szyku\l a.
\newblock Synchronizing Automata: Open Problems.
\newblock \emph{Electronic Proceedings in Theoretical Computer Science},
451:33--47, 2026.
\newblock DOI: \href{https://doi.org/10.4204/EPTCS.451.3}{10.4204/EPTCS.451.3}.
\newblock \url{https://arxiv.org/abs/2608.24245}.
\newblock Version~1, 25 August 2026; in particular Section~2.3,
Figure~3 and Conjecture~8. Consulted 20 September 2026.

\bibitem{MM2026}
Korea Superintelligence Labs; generated by Machina Mathematica.
\newblock The compress-with-another threshold of the series $A_k$: an upper
bound $4n/3$ for every $k$, exact values to $n=186$, and a one-letter correction.
\newblock \emph{Ideosphere}, version~1, source snapshot 7 September 2026.
\newblock \url{https://ideosphere.ai/papers/cwa-series}.
\newblock Archived manuscript:
\url{https://ideosphere.ai/papers/cwa-series-v1.pdf}.
\newblock AI-generated research manuscript; independent human review listed
as open. Consulted 20 September 2026. Its reported Lean build was not rerun
for the present article.

\bibitem{Zhang2026}
Enkai Zhang.
\newblock The triple rendezvous time of a synchronizing automaton can be
$\lfloor4n/3\rfloor$.
\newblock Preprint, arXiv:2609.19173v1, 14 September 2026.
\newblock \url{https://arxiv.org/abs/2609.19173}.
\newblock Consulted 20 September 2026. This paper concerns a different
family and is cited only to distinguish the neighboring invariant.

\end{thebibliography}
\end{document}
